% ============================================================
% Chapter IV: Experiments and Results
% ============================================================
\clearpage
\section{EXPERIMENTS AND RESULTS}

\subsection{Experimental Setup}

\subsubsection{Hardware Configuration}

Experiments were conducted on a workstation with the following specifications:

\begin{table}[htbp]
\centering
\caption{Hardware Configuration}
\label{tab:hardware}
\begin{tabular}{@{}ll@{}}
\toprule
\textbf{Component} & \textbf{Specification} \\
\midrule
CPU & Intel Core i7-13700K @ 3.4 GHz \\
GPU & NVIDIA RTX 4070 Ti SUPER 16 GB \\
RAM & 32 GB DDR5-5600 \\
Storage & 1 TB NVMe SSD \\
\bottomrule
\end{tabular}
\end{table}

\subsubsection{Software Environment}

\begin{table}[htbp]
\centering
\caption{Software Versions}
\label{tab:software}
\begin{tabular}{@{}ll@{}}
\toprule
\textbf{Software} & \textbf{Version} \\
\midrule
Operating System & Windows 11 (22631) \\
Unreal Engine & 5.4.4 \\
ASVSim & CosysAirSim v3.0.1 \\
Python & 3.10.11 \\
OpenCV & 4.8.1 \\
NumPy & 1.24.3 \\
\bottomrule
\end{tabular}
\end{table}

\subsubsection{Simulation Environment}

The LakeEnv environment provided with ASVSim was used for validation. This environment features:

\begin{itemize}
    \item Open water surface with physically simulated waves
    \item Distant shoreline with procedural terrain
    \item Dynamic lighting with day/night cycle
    \item Realistic water rendering with reflections and refraction
\end{itemize}

While not a polar environment, LakeEnv provides a suitable testbed for validating the data collection pipeline, with water surface dynamics similar to ice-free polar waters.

\subsection{Sensor Validation}

Prior to full-scale data collection, individual sensors were validated using the \texttt{2-verify\_sensors.py} script. The validation confirmed:

\begin{table}[htbp]
\centering
\caption{Sensor Validation Results}
\label{tab:sensor_validation}
\begin{tabular}{@{}llcl@{}}
\toprule
\textbf{Sensor} & \textbf{Test} & \textbf{Result} & \textbf{Notes} \\
\midrule
Connection & ASVSim client & \checkmark & Connection established \\
Vessel & Vessel1 spawn & \checkmark & Physics engine active \\
Front Camera RGB & Resolution 640$\times$480 & \checkmark & Color channels correct \\
Front Camera Depth & Resolution 640$\times$480 & \checkmark & Metric depth in meters \\
LiDAR & Point count & \checkmark & 8,192 points/scan \\
Camera Intrinsics & K matrix & \checkmark & $f_x$ = 466.0 calculated \\
\bottomrule
\end{tabular}
\end{table}

All six validation tests passed, confirming the sensor configuration was operational.

\subsection{Data Collection Experiments}

\subsubsection{Experiment 1: Baseline Circular Trajectory}

A 200-frame dataset was collected using the circular trajectory mode with optimized settings.

\textbf{Collection Parameters:}
\begin{itemize}
    \item Trajectory: Circular (constant thrust = 0.3, angle = 0.6)
    \item Frame count: 200
    \item Resolution: 640$\times$480
    \item Movement interval: 0.3 seconds
\end{itemize}

\textbf{Results:}

\begin{table}[htbp]
\centering
\caption{Data Collection Results (200 Frames)}
\label{tab:collection_results}
\begin{tabular}{@{}ll@{}}
\toprule
\textbf{Metric} & \textbf{Value} \\
\midrule
Total frames collected & 200 \\
Failed frames & 0 \\
Success rate & 100\% \\
Total time & 22.3 minutes \\
Average time per frame & 6.7 seconds \\
Output data size & 847 MB \\
\bottomrule
\end{tabular}
\end{table}

The collection achieved the target 200 frames with zero failures, validating the reliability of the optimized pipeline.

\subsubsection{Experiment 2: Performance Comparison}

The impact of rendering optimizations was quantified by comparing collection speeds before and after configuration changes.

\begin{table}[htbp]
\centering
\caption{Performance Comparison: Before and After Optimization}
\label{tab:performance}
\begin{tabular}{@{}lcc@{}}
\toprule
\textbf{Configuration} & \textbf{Time/Frame} & \textbf{Speedup} \\
\midrule
1280$\times$720 + Lumen & $\sim$55--60s & 1$\times$ (baseline) \\
640$\times$480 + No Lumen & $\sim$6--7s & \textbf{9$\times$} \\
\bottomrule
\end{tabular}
\end{table}

The 9$\times$ speed improvement demonstrates the critical importance of rendering configuration for practical data collection.

\subsubsection{Experiment 3: Trajectory Analysis}

The circular trajectory provides approximately 360$^\circ$ coverage of the surrounding environment. Analysis of vessel positions from the collected dataset shows:

\begin{itemize}
    \item Trajectory radius: approximately 15--20 meters
    \item Total path length: approximately 100--120 meters
    \item Average speed: 0.3--0.5 m/s
    \item Pose variation between frames: 1--2$^\circ$ rotation
\end{itemize}

This viewpoint distribution is suitable for 3D Gaussian Splatting training, providing sufficient angular coverage with adequate baseline between views.

\subsection{Data Quality Analysis}

\subsubsection{RGB Image Quality}

Sample RGB images from the collected dataset demonstrate:

\begin{itemize}
    \item \textbf{Color accuracy:} Water appears in natural blue tones after RGB-to-BGR conversion
    \item \textbf{Resolution adequacy:} 640$\times$480 captures sufficient detail for 3DGS training
    \item \textbf{Edge quality:} Removal of vignette and chromatic aberration effects produces clean image boundaries
    \item \textbf{Vessel exclusion:} Optimized camera position (X=3.5m, Z=-2.5m) successfully removes vessel structure from field of view
\end{itemize}

Figure \ref{fig:sample_rgb} shows representative frames from the dataset.

\begin{figure}[htbp]
\centering
% Placeholder for sample RGB images
\fbox{\parbox{0.45\textwidth}{\centering
\vspace{1.5cm}
Frame 0000\\
(RGB Image)
\vspace{1.5cm}
}}
\hfill
\fbox{\parbox{0.45\textwidth}{\centering
\vspace{1.5cm}
Frame 0100\\
(RGB Image)
\vspace{1.5cm}
}}
\caption{Sample RGB frames from the collected dataset showing water surface and shoreline.}
\label{fig:sample_rgb}
\end{figure}

\subsubsection{Depth Map Analysis}

Depth maps were analyzed for coverage and accuracy:

\begin{table}[htbp]
\centering
\caption{Depth Map Statistics}
\label{tab:depth_stats}
\begin{tabular}{@{}ll@{}}
\toprule
\textbf{Metric} & \textbf{Value} \\
\midrule
Valid depth pixels (median) & 82\% \\
Minimum depth & 1.2 m \\
Maximum depth & 65504 m (infinity) \\
Mean valid depth & 45.3 m \\
Depth standard deviation & 28.7 m \\
\bottomrule
\end{tabular}
\end{table}

Invalid depth pixels (approximately 18\%) correspond primarily to sky regions, which is expected and acceptable for 3DGS training. The depth distribution shows good coverage of the scene with values concentrated in the 10--80 meter range.

\subsubsection{LiDAR Point Cloud Quality}

LiDAR scans consistently produced the expected 8,192 points across all frames:

\begin{itemize}
    \item Point count consistency: 8,192 points (100\% of scans)
    \item Point distribution: Uniform in azimuth, structured in elevation
    \item Range coverage: 0--100 meters as configured
    \item Frame-to-frame consistency: No dropped scans
\end{itemize}

\subsubsection{Pose Continuity}

Camera pose data was analyzed for continuity and smoothness:

\begin{itemize}
    \item Position jumps: No anomalous position discontinuities detected
    \item Orientation smoothness: Quaternion interpolation shows smooth rotation
    \item Velocity consistency: Speed values match expected vessel dynamics
\end{itemize}

\subsection{Dataset Validation}

Automated validation using \texttt{validate\_dataset.py} confirmed:

\begin{table}[htbp]
\centering
\caption{Dataset Validation Summary}
\label{tab:validation}
\begin{tabular}{@{}lll@{}}
\toprule
\textbf{Check} & \textbf{Status} & \textbf{Details} \\
\midrule
File count consistency & PASS & 200 files in each modality \\
Image dimensions & PASS & All 640$\times$480$\times$3 \\
Depth shape & PASS & All 640$\times$480 \\
Depth valid pixels & PASS & Mean 82\% valid \\
LiDAR point count & PASS & 8,192 points/frame \\
Pose file completeness & PASS & 200 poses \\
Pose continuity & PASS & No large jumps \\
\bottomrule
\end{tabular}
\end{table}

All validation checks passed, confirming the dataset is ready for downstream processing.

\subsection{COLMAP Format Verification}

The \texttt{transforms.json} file generated for 3DGS training was verified:

\begin{lstlisting}[language={}, caption={Excerpt from transforms.json}]
{
  "camera_model": "OPENCV",
  "fl_x": 466.0, "fl_y": 466.0,
  "cx": 320.0, "cy": 240.0,
  "w": 640, "h": 480,
  "frames": [
    {
      "file_path": "rgb/0000.png",
      "transform_matrix": [...]
    },
    ...
  ]
}
\end{lstlisting}

The file contains 200 frame entries with properly formatted camera intrinsics and transformation matrices, compatible with standard 3D Gaussian Splatting training scripts.

\subsection{Discussion of Results}

The experimental results demonstrate several key achievements:

\textbf{Efficiency:} The 9$\times$ performance improvement through rendering optimization makes large-scale data collection practical. A 200-frame dataset that would require 3 hours at original settings now completes in 22 minutes.

\textbf{Reliability:} Zero frame failures across 200 frames indicates the error handling and retry mechanisms are effective.

\textbf{Quality:} Validation confirms the collected data meets specifications for 3DGS training---synchronized multi-modal data with accurate poses.

\textbf{Completeness:} The pipeline produces all necessary outputs including RGB, depth, LiDAR, poses, and COLMAP-formatted transforms.

These results validate the data collection pipeline as a solid foundation for subsequent phases including 3DGS training and path planning development.
